\documentclass[11pt]{article}

\usepackage[margin=1in]{geometry}
\usepackage{amsmath,amssymb,amsthm}
\usepackage{bm}
\usepackage{mathtools}
\usepackage{booktabs}
\usepackage{enumitem}
\usepackage{array}
\usepackage{hyperref}

\newcommand{\R}{\mathbb{R}}
\newcommand{\E}{\mathbb{E}}
\newcommand{\dom}{\Omega}
\newcommand{\bdry}{\partial\Omega}
\newcommand{\Gtrue}{\mathcal{G}^{\star}}
\newcommand{\Gprior}{\widetilde{\mathcal{G}}}
\newcommand{\Ghat}{\widehat{\mathcal{G}}}
\newcommand{\Ntrue}{\mathcal{N}^{\star}}
\newcommand{\Nass}{\widetilde{\mathcal{N}}}
\newcommand{\Rcorr}{\mathcal{R}}
\newcommand{\lap}{\Delta}
\newcommand{\dt}{\partial_t}
\newcommand{\dx}{\partial_x}
\newcommand{\dxx}{\partial_{xx}}
\newcommand{\dxxx}{\partial_{xxx}}
\newcommand{\norm}[1]{\lVert #1 \rVert}

\title{Correcting Model Misspecification in Neural-Operator Physics Priors}
\author{Shreeyam Bangera}
\date{\today}

\begin{document}
\maketitle

\begin{abstract}
A neural operator is trained without solution data, using only the residual of an assumed PDE
and its boundary or initial conditions; this yields a \emph{physics prior} that maps a problem
input to an approximate solution. When the assumed PDE differs from the process that generated
the data---parametrically or structurally---the prior carries a systematic, data-irreducible
error. We formalize this model-misspecification error, state an architecture-agnostic scheme
that corrects it with a small labeled set by learning an additive correction on top of the
frozen prior, list the models and hyperparameters used throughout, and describe a benchmark of
five PDEs, each with four graded misspecifications spanning wrong coefficients through wrong
operator structure.
\end{abstract}

\section{The physics prior}
\label{sec:setup}

Let $\dom\subset\R^{d}$ be a bounded domain. A problem instance is specified by an input
function $v\in\mathcal{V}$ (a forcing term for the steady problems, an initial condition for the
time-dependent ones), drawn from a probability measure $\mu$ on $\mathcal{V}$. The true process
is governed by a differential or integro-differential operator $\Ntrue$ with a boundary/initial
constraint $\mathcal{C}$,
\begin{equation}
\label{eq:true}
\Ntrue\!\left(u, v\right)(z) = 0,\quad z\in\dom,
\qquad
\mathcal{C}\!\left(u, v\right)(z) = 0,\quad z\in\Gamma,
\end{equation}
where $\Gamma$ is the boundary $\bdry$ (Dirichlet data) or the initial slice $\{t=0\}$.
Equation~\eqref{eq:true} defines the true solution map $\Gtrue:v\mapsto u$.

We train a neural operator $\mathcal{G}_\theta$ \emph{without any solution labels}. Given an
assumed model $\Nass$ and the same constraint $\mathcal{C}$, we minimize the expected squared
residual over interior collocation points $X_c\subset\dom$ and the constraint over points
$X_\Gamma\subset\Gamma$,
\begin{equation}
\label{eq:physloss}
\mathcal{L}_{\mathrm{phys}}(\theta)
= \E_{v\sim\mu}\!\left[
\frac{1}{|X_c|}\sum_{z\in X_c}\bigl|\Nass\!\left(\mathcal{G}_\theta(v), v\right)(z)\bigr|^{2}
\right]
+ \lambda\,
\E_{v\sim\mu}\!\left[
\frac{1}{|X_\Gamma|}\sum_{z\in X_\Gamma}\bigl|\mathcal{C}\!\left(\mathcal{G}_\theta(v), v\right)(z)\bigr|^{2}
\right].
\end{equation}
Collocation points are resampled each step and no term in \eqref{eq:physloss} uses the target
$u=\Gtrue(v)$. Write the minimizer as the \emph{prior operator}
$\Gprior = \arg\min_\theta \mathcal{L}_{\mathrm{phys}}(\theta)$. If $\Nass=\Ntrue$, then
$\Gprior$ recovers the true solution map up to optimization and approximation error, and
$\mathcal{L}_{\mathrm{phys}}$ is a consistent surrogate for the supervised loss.

\section{Model misspecification}
\label{sec:misspec}

\subsection{What we are trying to do}
\label{sec:theory}
In practice the assumed model is wrong: $\Nass\neq\Ntrue$, so the prior solves a different
problem than the one that produced the data. Define the \emph{model-form error} of the prior on
input $v$,
\begin{equation}
\label{eq:bias}
\varepsilon(v) \;=\; \Gprior(v) - \Gtrue(v).
\end{equation}
Two properties distinguish $\varepsilon$ from ordinary estimation error.
\begin{enumerate}[leftmargin=1.4em,itemsep=2pt]
\item \textbf{It is systematic.} $\varepsilon$ is fixed by the gap $\Nass-\Ntrue$, not by finite
sampling. Adding collocation points, widening the network, or training longer drives
$\mathcal{G}_\theta$ toward $\Gprior$ and therefore \emph{toward}, not away from, the biased
solution.
\item \textbf{It is structured.} Because $\Gprior(v)$ still satisfies the assumed physics and the
correct constraints, $\varepsilon(v)$ inherits regularity from both operators and is typically
smoother and smaller in amplitude than $u$ itself.
\end{enumerate}
We separate two regimes. \emph{Parametric} misspecification keeps the operator form and perturbs
a coefficient (a wrong diffusivity or wavenumber). \emph{Structural} misspecification changes
the operator class itself: linearizing a nonlinear term, replacing a nonlocal operator by a local
one, adding a term absent from the true model, or turning a constant coefficient into a spatially
varying field. The aim is to quantify, across both regimes and at graded severity, how large
$\varepsilon$ is and how cheaply it can be removed with a small amount of labeled data
(Section~\ref{sec:correct}).

\subsection{Physics-prior hyperparameters}
\label{sec:prior-hparams}
The prior $\Gprior$ is a DeepONet in every problem: a branch network encodes the input $v$, a
trunk network encodes the query coordinate, and their $P$-dimensional codes are combined by an
inner product. For the 1D and space--time problems (P1--P4) the branch and trunk are
$\tanh$-MLPs; the trunk of P2--P4 takes truncated Fourier features of the coordinate (order $M$).
For the 2D problem (P5) the branch is a convolutional encoder of the source image and the trunk
takes random Fourier features (scales $\{6,12,20\}$, $K=64$). Derivatives entering the residual
$\Nass$ are taken by automatic differentiation of the trunk; the nonlocal (fractional) operator
of P4 is evaluated spectrally. Table~\ref{tab:prior} lists the sizes and the physics-training
schedule. All training uses Adam at learning rate $10^{-3}$, seed $0$, and early stopping on the
validation physics loss with relative tolerance $10^{-4}$.

\begin{table}[h]
\centering
\footnotesize
\setlength{\tabcolsep}{5pt}
\begin{tabular}{@{}lllcccr@{}}
\toprule
& Branch & Trunk & $P$ & Epochs/pat. & Colloc.\ $N_c$ & Params \\
\midrule
P1 & MLP $[128,128]$ & MLP $[64,64,64]$ & 64 & 1500 / 60 & 101 & $50{,}433$ \\
P2 & MLP $[256,256]$ & Fourier$_{16}$ + MLP $[128,128,128]$ & 128 & 800 / 50 & 3000 & $218{,}369$ \\
P3 & MLP $[256,256]$ & Fourier$_{16}$ + MLP $[128,128,128]$ & 128 & 800 / 50 & 3000 & $218{,}369$ \\
P4 & MLP $[256,256]$ & Fourier$_{20}$ + MLP $[128,128,128]$ & 128 & 800 / 50 & 3000 & $284{,}929$ \\
P5 & CNN $[16,32,64,64]$ & RFF$_{384}$ + MLP $[256,256,256]$ & 128 & 800 / 50 & 2500 & $1{,}404{,}993$ \\
\bottomrule
\end{tabular}
\caption{Physics-prior architecture, training schedule, and total parameter count. Fourier$_{M}$
denotes trunk Fourier features of order $M$; RFF$_{384}$ the $384$-dimensional random Fourier
feature map; the P5 branch is a four-layer stride-2 CNN with the listed channel widths. Steps per
epoch: $8$ (P1, P5), $16$ (P2--P4). Physics batch (inputs per step): $256$ (P1), $32$ (P2--P3),
$64$ (P4), $8$ (P5). Boundary/initial data: Dirichlet penalty $\lambda{=}10$ (P1), initial-condition
penalty $\lambda{=}20$ (P2--P4); P5 enforces homogeneous Dirichlet data exactly by multiplying the
output by the bump $x(1-x)y(1-y)$, so it needs no penalty. Optimizer Adam, learning rate
$10^{-3}$, seed $0$, early stopping on the validation physics loss (relative tolerance $10^{-4}$).}
\label{tab:prior}
\end{table}

\section{Correcting the misspecified prior}
\label{sec:correct}

We assume a \emph{small} labeled set $\mathcal{D}=\{(v_i,u_i)\}_{i=1}^{N}$ with
$u_i=\Gtrue(v_i)$, alongside the unlabeled distribution $\mu$ used for the prior. The prior is
treated as a fixed, informative guess and is not retrained.

\paragraph{Additive correction.}
Freeze $\Gprior$ and introduce a second neural operator, the correction
$\Rcorr_\phi:\mathcal{V}\times\mathcal{U}\to\mathcal{U}$, that takes the input and the prior
prediction and outputs a residual field. The corrected estimator is
\begin{equation}
\label{eq:combined}
\Ghat(v) \;=\; \Gprior(v) \;+\; \Rcorr_\phi\!\left(v,\ \Gprior(v)\right),
\end{equation}
trained by supervised regression,
\begin{equation}
\label{eq:dataloss}
\mathcal{L}_{\mathrm{data}}(\phi)
= \frac{1}{N}\sum_{i=1}^{N}
\norm{\,\Gprior(v_i) + \Rcorr_\phi\!\left(v_i,\Gprior(v_i)\right) - u_i\,}_{L^2(\dom)}^{2}.
\end{equation}
At the optimum $\Rcorr_\phi(v,\Gprior(v))\approx u-\Gprior(v)=-\varepsilon(v)$: the correction
learns the model-form error itself. When the prior captures the dominant structure of the
solution, $\varepsilon$ is a low-amplitude, smooth field, so $\Rcorr_\phi$ has an easier target
than $\Gtrue$ and is learnable from few labels. The scheme is agnostic to the operator
architecture and to the specific misspecification; the prior enters only through the frozen map
$v\mapsto\Gprior(v)$.

\paragraph{Baselines.}
To isolate the value of the prior we compare \eqref{eq:combined} against a \emph{pure-data}
operator $\mathcal{H}_\psi$ trained from scratch on the same labels,
$\min_\psi \tfrac{1}{N}\sum_i \norm{\mathcal{H}_\psi(v_i)-u_i}_{L^2}^2$, with no prior. For each
PDE and level we report the prior's held-out physics residual, the prior's own relative $L^2$
error $\norm{\Gprior(v)-u}/\norm{u}$, the pure-data operators, and the combined estimator
\eqref{eq:combined}. The contrast, read against the prior error, measures how much a
\emph{wrong} model still helps once corrected.

\subsection{Estimators and hyperparameters}
\label{sec:corr-hparams}
Both the correction $\Rcorr_\phi$ and the pure-data baseline $\mathcal{H}_\psi$ are instantiated
as a DeepONet and, separately, as a Fourier neural operator (FNO), so each problem yields four
data-trained models: pure-data DeepONet, pure-data FNO, combined DeepONet, combined FNO
(P5 uses FNO only). The combined variants add one input to the same backbone---the DeepONet
trunk takes the prior value as an extra feature, the FNO takes the prior field as an extra
channel (for P1 the combined DeepONet branch additionally ingests the prior field, doubling its
input width). The data DeepONet has the same trunk/coordinate encoding as the prior but a smaller
capacity; the FNO is a standard spectral-convolution stack. Table~\ref{tab:corr} lists sizes and
the supervised schedule. All training uses Adam at learning rate $10^{-3}$, seed $0$, and early
stopping on the validation loss with relative tolerance $5\times10^{-3}$.

\begin{table}[h]
\centering
\footnotesize
\setlength{\tabcolsep}{5pt}
\begin{tabular}{@{}llrlrcc@{}}
\toprule
& Data DeepONet ($P{=}64$) & DON par. & Data FNO ($w,m,L$) & FNO par. & Ep./pat. & Batch (DON/FNO) \\
\midrule
P1 & $[128,128]/[64,64,64]$ & $50{,}433$ & FNO1d $24,\,16,\,2$ & $23{,}033$ & 300 / 20 & 256 / 256 \\
P2 & $[128,128]/[64,64,64]$ & $72{,}321$ & FNO2d $12,\,(8,8),\,2$ & $39{,}017$ & 200 / 20 & 64 / 16 \\
P3 & $[128,128]/[64,64,64]$ & $72{,}321$ & FNO2d $12,\,(8,8),\,2$ & $39{,}017$ & 200 / 20 & 64 / 16 \\
P4 & $[128,128]/[64,64,64]$ & $105{,}601$ & FNO2d $12,\,(8,8),\,2$ & $39{,}017$ & 200 / 20 & 64 / 8 \\
P5 & --- (FNO only) & --- & FNO2d $32,\,(24,24),\,4$ & $4{,}727{,}297$ & 200 / 20 & --- / 8 \\
\bottomrule
\end{tabular}
\caption{Correction and baseline architectures, parameter counts, and supervised training.
Parameters are for the pure-data model; the combined variant adds one input (the prior value or
field): the DeepONet trunk gains one feature ($+64$ params), except for P1 where the combined
DeepONet also doubles its branch input to $202$, giving $63{,}361$; the FNO gains one lifting
channel ($+12$ to $+32$ params). Modes $m$ are retained Fourier modes per axis, $w$ the channel
width, $L$ the number of spectral layers. The DeepONet loss samples $4000$ grid points per input
(P2--P4), $4096$ (P5, unused there), or the full grid (P1); the FNO predicts the whole field.
Optimizer Adam, learning rate $10^{-3}$, seed $0$, early stopping on the validation loss (relative
tolerance $5\times10^{-3}$).}
\label{tab:corr}
\end{table}

\section{Benchmark problems and their misspecifications}
\label{sec:data}

Each dataset is a set of input/solution pairs $(v,u)$ generated from a fixed true model, split
$2048/200/500$ into train/validation/test. Inputs are Gaussian random fields (P2--P5) or a
manufactured random combination of modes (P1); every solution is computed in closed form---P1
analytically, P3/P4 by the diagonal Fourier propagator, P5 by the diagonal Dirichlet sine
transform---or spectrally converged (P2, pseudospectral RK4 at $512$ modes, coarsened to $256$).
The datasets therefore carry no discretization error, so the discrepancies the benchmark measures
are model-form error alone, not solver error. Table~\ref{tab:gen} collects domains, grids, and
input laws. The prior for a given level minimizes \eqref{eq:physloss} with the \emph{assumed}
residual below; the four levels \textsf{low}, \textsf{med}, \textsf{high}, \textsf{max} are ordered
by increasing departure from the true operator. Here $v$ is the input and $u$ the solution;
$C_0=\big(\E_x v^2\big)^{1/2}$ is the root-mean-square input amplitude.

\begin{table}[h]
\centering
\footnotesize
\setlength{\tabcolsep}{5pt}
\begin{tabular}{@{}lllll@{}}
\toprule
& Domain & Grid & Input law & Solution \\
\midrule
P1 & $x\in[-1,1]$ & $101$ & manufactured $\tfrac{x^2-1}{10}\sum_{i=1}^{5}(w^s_i\sin i\pi x+w^c_i\cos i\pi x),\ w\!\sim\!U(0,1)$ & analytic \\
P2 & $x\in[0,1)$ per., $t\in[0,0.5]$ & $256\times101$ & GRF $\mathcal{N}\!\big(0,\,25^2(-\lap+25I)^{-2}\big)$ & pseudospectral RK4 \\
P3 & $x\in[0,1)$ per., $t\in[0,0.5]$ & $256\times101$ & GRF $\mathcal{N}\!\big(0,\,25^2(-\lap+25I)^{-2}\big)$ & exact Fourier \\
P4 & $x\in[0,1)$ per., $t\in[0,1]$ & $512\times101$ & GRF $\mathcal{N}\!\big(0,\,5^2(-\lap+25I)^{-1.25}\big)$ & exact Fourier \\
P5 & $(0,1)^2$, interior & $128\times128$ & sine-basis GRF, $c_{mn}\!\sim\!\mathcal{N}\!\big(0,\,836.5^2(\lambda_{mn}+25)^{-1.5}\big)$ & exact sine transform \\
\bottomrule
\end{tabular}
\caption{Data generation. All splits are $2048/200/500$ (train/val/test). ``Grid'' is $N_x$ (P1),
$N_x\times N_t$ (P2--P4), or the interior $N\times N$ (P5). $\lambda_{mn}=\pi^2(m^2+n^2)$ are the
Dirichlet eigenvalues. The P4 input measure uses exponent $1.25$ (a heavier spectral tail than the
$2$ of P2/P3), so its nonlocal-versus-local discrepancy is not masked by a low-wavenumber input.}
\label{tab:gen}
\end{table}

\subsection*{P1 --- Nonlinear diffusion--reaction (1D, steady)}
Domain $\dom=(-1,1)$ on a $101$-point grid, input forcing $v(x)$, homogeneous Dirichlet data
$u(-1)=u(1)=0$. True model:
\begin{equation}
D\,\dxx u - K_0\, e^{-u} u - v = 0,
\qquad D=0.1,\ K_0=0.5 .
\end{equation}
Assumed residuals driven to zero by the prior:
\begin{align*}
\textsf{low:}\ & 0.11\,\dxx u - 0.5\,e^{-u} u - v
&&\text{(diffusivity $+10\%$)}\\
\textsf{med:}\ & 0.15\,\dxx u - 0.7\,e^{-u} u - v
&&\text{(diffusivity and reaction rate off)}\\
\textsf{high:}\ & 0.1\,\dxx u - 0.5\, u - v
&&\text{(reaction linearized, $e^{-u}u\!\to\!u$)}\\
\textsf{max:}\ & \dx\!\big(D(x)\,\dx u\big) - 0.5\,e^{-u} u - v,\quad D(x)=0.1\,(1+0.5\cos 2\pi x)
&&\text{(constant $D\!\to\!$ heterogeneous)}
\end{align*}

\subsection*{P2 --- Viscous Burgers (1+1D)}
Space--time domain with periodic $x$, initial condition $u(x,0)=v(x)$.
True model:
\begin{equation}
\dt u + u\,\dx u - \nu\,\dxx u = 0, \qquad \nu = 0.01 .
\end{equation}
Assumed residuals:
\begin{align*}
\textsf{low:}\ & \dt u + u\,\dx u - 0.013\,\dxx u &&\text{(viscosity $+30\%$)}\\
\textsf{med:}\ & \dt u + u\,\dx u - 0.03\,\dxx u &&\text{(viscosity $3\times$)}\\
\textsf{high:}\ & \dt u + C_0\,\dx u - 0.01\,\dxx u &&\text{(advection linearized, $u\dx u\!\to\!C_0\dx u$)}\\
\textsf{max:}\ & \dt u + u\,\dx u + \delta\,\dxxx u,\quad \delta=10^{-3} &&\text{(viscosity dropped, KdV dispersion added)}
\end{align*}

\subsection*{P3 --- Advection--diffusion (1+1D)}
Periodic $x$, initial condition $u(x,0)=v(x)$.
True model:
\begin{equation}
\dt u + c\,\dx u - \nu\,\dxx u = 0, \qquad c=1,\ \nu=10^{-3}.
\end{equation}
Assumed residuals:
\begin{align*}
\textsf{low:}\ & \dt u + 1.05\,\dx u - \nu\,\dxx u &&\text{(wave speed $+5\%$)}\\
\textsf{med:}\ & \dt u + 1.3\,\dx u - \nu\,\dxx u &&\text{(wave speed $+30\%$)}\\
\textsf{high:}\ & \dt u + c(x)\,\dx u - \nu\,\dxx u,\quad c(x)=1+0.3\sin 2\pi x &&\text{(constant $c\!\to\!$ variable)}\\
\textsf{max:}\ & \dt u + (1+u)\,\dx u - \nu\,\dxx u &&\text{(linear $\to$ nonlinear advection)}
\end{align*}

\subsection*{P4 --- Fractional diffusion (1+1D)}
Periodic $x$, initial condition $u(x,0)=v(x)$. The fractional Laplacian $(-\lap)^{s}$ acts in
Fourier space through the symbol $|2\pi k|^{2s}$.
True model:
\begin{equation}
\dt u + \kappa\,(-\lap)^{s} u = 0, \qquad \kappa=0.03,\ s=0.75 .
\end{equation}
Assumed residuals:
\begin{align*}
\textsf{low:}\ & \dt u + 0.033\,(-\lap)^{0.75} u &&\text{(coefficient $+10\%$)}\\
\textsf{med:}\ & \dt u + 0.03\,(-\lap)^{0.65} u &&\text{(wrong fractional order $s$)}\\
\textsf{high:}\ & \dt u - 0.03\,\dxx u &&\text{(nonlocal $\to$ local diffusion)}\\
\textsf{max:}\ & \dt u - \dx\!\big(a(x)\,\dx u\big),\quad a(x)=0.03\,(1+0.5\cos 2\pi x) &&\text{(nonlocal $\to$ heterogeneous local)}
\end{align*}

\subsection*{P5 --- Helmholtz (2D)}
Domain $\dom=(0,1)^2$ on a $128\times128$ interior grid, source $v(x,y)$, homogeneous Dirichlet
data $u|_{\bdry}=0$. True model:
\begin{equation}
-\lap u - \kappa^{2} u - v = 0, \qquad \kappa^{2}=3765.25\ \ (\kappa=61.362).
\end{equation}
The value $\kappa^2=3765.25$ is the midpoint of the widest gap in the Dirichlet spectrum below
$9000$ (between eigenvalues $3720.8$ and $3809.7$), so the true operator is well conditioned. The
\textsf{med} shift to $3805$ sits only $\approx 4.7$ from the eigenvalue $3809.7$, driving the
resolvent toward a pole: a small parameter error produces an $O(1)$ change in $u$, so here
misspecification magnitude and discrepancy size are not proportional.
Assumed residuals:
\begin{align*}
\textsf{low:}\ & -\lap u - 3780\,u - v &&\text{(wavenumber $\kappa^2$ off)}\\
\textsf{med:}\ & -\lap u - 3805\,u - v &&\text{(wavenumber $\kappa^2$ off more)}\\
\textsf{high:}\ & -\nabla\!\cdot\!\big(a\,\nabla u\big) - \kappa^{2} u - v,\quad a=1+0.2\cos 2\pi x\,\cos 2\pi y
&&\text{(homogeneous $\to$ smooth heterogeneous medium)}\\
\textsf{max:}\ & -\nabla\!\cdot\!\big(a\,\nabla u\big) - \kappa^{2} u - v,\quad a=1+1.5\big(1+\tanh\tfrac{0.25-r}{0.02}\big)
&&\text{(high-contrast radial scatterer)}
\end{align*}
with $r=\sqrt{(x-\tfrac12)^2+(y-\tfrac12)^2}$.

\subsection*{Summary}
Table~\ref{tab:taxonomy} classifies each level as parametric (wrong coefficient, correct form) or
structural (wrong operator class). Severity grows left to right, and every problem reaches a
structural misspecification at \textsf{max}.

\begin{table}[h]
\centering
\begin{tabular}{@{}llllll@{}}
\toprule
& \textsf{low} & \textsf{med} & \textsf{high} & \textsf{max}\\
\midrule
P1 Diffusion--reaction & param.\ ($D$) & param.\ ($D,K_0$) & struct.\ (linearized) & struct.\ (heterog.\ $D$)\\
P2 Burgers & param.\ ($\nu$) & param.\ ($\nu$) & struct.\ (linearized) & struct.\ (dispersive)\\
P3 Advection--diffusion & param.\ ($c$) & param.\ ($c$) & struct.\ (variable $c$) & struct.\ (nonlinear)\\
P4 Fractional diffusion & param.\ ($\kappa$) & struct.\ (order $s$) & struct.\ (local) & struct.\ (heterog.\ local)\\
P5 Helmholtz & param.\ ($\kappa^2$) & param.\ ($\kappa^2$) & struct.\ (heterog.) & struct.\ (high contrast)\\
\bottomrule
\end{tabular}
\caption{Misspecification type per problem and level.}
\label{tab:taxonomy}
\end{table}

\end{document}
